\section{Quarter-shifted paired hypergeometric identities}
\label{sec:paired-case}\record{Q04}\record{Q05}
Leibtag et al.~\cite[Section 6.2]{Leibtag} represent the same
approximation coefficients in two ways: a terminating hypergeometric
sum and minors of a dual coefficient matrix. Different presentations
can expose different prime divisors. The experiment here deforms that
pair of presentations to the quarter arguments $1/4$ and $3/4$, then
subtracts to isolate $\theta$. Shifting only the final harmonic sums
does not preserve the identity. The kernels, their coefficients and
the normalizing scalar all have to change together.

This is the most computationally substantial case study. We give the
finite definitions, the exact remainder formulas, and the mechanism
of the paired identities. The accompanying directory
\path{checks/paired/} contains the rational matrix data and a standalone
exact replay of the ten polynomial certificates. Those data are part
of this case study, not a reference to an unspecified research log.
The mathematical argument explains what the certificates prove and
why their boundary terms vanish.

\subsection{Forward coefficients and an exact remainder}
Let $\alpha=-3/4$ or $-1/4$, and put $a=1+\alpha$. For $n\geq0$ set
\begin{equation}\label{eq:paired-B}
 B_n^\alpha(k)=\prod_{j\in\{12,13,14\}}
                 \binom{k+jn+\alpha}{10n}
                     \binom{15n}{n+k}\binom{15n}{k}.
\end{equation}
For differentiation this means the full Gamma interpolation, with
$\binom{x}{m}=x(x-1)\cdots(x-m+1)/m!$ for integer $m\geq0$ and
$\binom{15n}{z}=\Gamma(15n+1)/[\Gamma(z+1)\Gamma(15n-z+1)]$.
In particular, $B_n^\alpha(k)$ can be zero at an integer while its
derivative is nonzero. Put
\[
 H_m^{(s)}(a)=\sum_{j=0}^{m-1}(j+a)^{-s},
\]
and define the rational numbers
\begin{align}
 Q_n^a&=\sum_{k=0}^{14n}B_n^\alpha(k),\notag\\
 P_{s,n}^a&=\sum_{k=-n}^{15n}\left(
 B_n^\alpha(k)H_{4n+k}^{(s)}(a)
 -\frac{(B_n^\alpha)'(k)}{s-1}H_{4n+k}^{(s-1)}(a)\right),
 \quad s=3,4.\label{eq:paired-forward}
\end{align}
The summation range in the second line includes both simple-zero tails.
For reproduction without Gamma differentiation, the derivative in the
interior $0\leq k\leq14n$ is $B_n^\alpha(k)$ times
\[
 \sum_{j=12}^{14}\sum_{i=0}^{10n-1}\frac1{k+jn+\alpha-i}
 +H_{14n-k}^{(1)}(1)-H_{n+k}^{(1)}(1)
 +H_{15n-k}^{(1)}(1)-H_k^{(1)}(1).
\]
For $14n<k\leq15n$, replace the last two binomials in
\eqref{eq:paired-B} by
$\binom{15n}{k}(-1)^k/[(n+k)\binom{n+k-1}{15n}]$;
for $k=-i$, $1\leq i\leq n$, replace them by
$\binom{15n}{n-i}(-1)^{i-1}/[(15n+i)\binom{15n+i-1}{i-1}]$.
These are exactly their simple-zero derivatives.

Write $(z)_m=z(z+1)\cdots(z+m-1)$ and define
\begin{align}
 \mathcal H_n(s)&=\frac{(-1)^n(15n)!^2}{(10n)!^3}
 \frac{\prod_{j=12}^{14}(s-jn-\alpha)_{10n}}
                         {(s)_{14n+1}(s-n)_{16n+1}},\notag\\
 T_n(z)&=\mathcal H_n(z+4n).
 \label{eq:paired-T}
\end{align}
Gamma reflection gives
$B_n^\alpha(x)=(\sin(\pi x)/\pi)^2T_n(-x-4n)$.
Taking the double and simple Laurent coefficients therefore gives
\begin{equation}\label{eq:paired-partial}
 T_n(z)=\sum_{k=-n}^{15n}\left(
 \frac{B_n^\alpha(k)}{(z+4n+k)^2}
 -\frac{(B_n^\alpha)'(k)}{z+4n+k}\right).
\end{equation}
Both sides are proper rational functions with the same principal parts;
the difference is a polynomial vanishing at infinity. Also
$\sum_k(B_n^\alpha)'(k)=0$ because $T_n=O(z^{-2})$.
Differentiate and sum~\eqref{eq:paired-partial} at $z=m+\alpha$,
$m\geq1$. The elementary identity
$\sum_{m\geq1}(m+\alpha+h)^{-s}=\zeta(s,a)-H_h^{(s)}(a)$
then proves
\begin{equation}\label{eq:paired-remainders}
 \begin{split}
 Q_n^a\zeta(3,a)-P_{3,n}^a
      &=-\frac12\sum_{m\geq1}T_n'(m+\alpha),\\
 Q_n^a\zeta(4,a)-P_{4,n}^a
      &=\frac16\sum_{m\geq1}T_n''(m+\alpha).
 \end{split}
\end{equation}
All differentiated tails converge absolutely. Their first $9n$ and
$8n$ sampled terms, respectively, vanish by the multiplicities of the
numerator roots. These formulas already define useful exact starting
points independently of the paired representation.

\subsection{Dual coefficient rows and the paired identities}
Define a different rational function
\begin{equation}\label{eq:paired-R}
 \mathcal R_n(s)=
 \frac{(-1)^n(1+2n-s)_{14n}(1+n-s)_{16n}}
 {(s+\alpha)_{11n+1}(s+n+\alpha)_{11n+1}(s+2n+\alpha)_{8n+2}}.
\end{equation}
For $k=0,1,2$, take the partial fractions
\[
 (-s)^k\mathcal R_n(s)=
       \sum_{j\geq0}\sum_{i=1}^3\frac{c_{j,i}^{(k)}}{(s+\alpha+j)^i}.
\]
Absent poles and orders have coefficient zero. Define three rational
row vectors, whose coordinates are indexed by $k$, by
\begin{equation}\label{eq:paired-rows}
 A_k=3\sum_jc_{j,3}^{(k)},\qquad
 B_k=2\sum_jc_{j,2}^{(k)},\qquad
 D_k=-\sum_{j,i}i\,c_{j,i}^{(k)}H_{2n+j}^{(i+1)}(a).
\end{equation}
The simple coefficients sum to zero, since the rational function is
$O(s^{-2})$ or smaller. Thus differentiation and summation yield the
vector identity
\begin{equation}\label{eq:paired-state}
 G_n=A_n\zeta(4,a)+B_n\zeta(3,a)+D_n
 =\left(-\sum_{m>2n}\left.\frac{d}{ds}
                  [(-s)^k\mathcal R_n(s)]\right|_{s=m}\right)_{k=0}^2.
\end{equation}
This proves the formal absence of $\zeta(2,a)$ without a numerical
linear-independence assumption. At $n=0$ the rows are
\[
 A_0=3(1,\alpha,\alpha^2),\quad
 B_0=-2(1,1+\alpha,\alpha^2+2\alpha),\quad
 D_0=(a^{-2},a^{-1},1).
\]

Put $u_n=(1,10n+a,(10n+a)^2)$ and
\begin{equation}\label{eq:paired-lambda}
 \Lambda_n=
 \frac{(8n)!(9n)!^2(11n)!^2(12n)!}{6(14n)!(16n)!}
 \frac{(a)_{11n}(a)_{12n}^2(a)_{13n}^2(a)_{14n}}
 {(a)_{16n}(a)_{17n}^2(a)_{18n}^2(a)_{19n}}.
\end{equation}
The complete paired formulas are
\begin{equation}\label{eq:paired-minors}
 Q_n^a=\Lambda_n\det\begin{bmatrix}B_n\\A_n\\u_n\end{bmatrix},\qquad
 P_{3,n}^a=-\Lambda_n\det\begin{bmatrix}D_n\\A_n\\u_n\end{bmatrix},\qquad
 P_{4,n}^a=\Lambda_n\det\begin{bmatrix}D_n\\B_n\\u_n\end{bmatrix}.
\end{equation}
The next subsection details the finite certificate proof of these
all-index identities. It is longer than merely checking several $n$,
because the derivative tails and the separate coefficient rows need
their own recurrence justification.

\subsection{How to verify the all-index identities}\label{subsec:paired-proof}
The data files specify rational matrices $M(n)$ of size $4$, $Q(n)$
of size $3$, a $4\times3$ matrix $K(n)$, and rank-three companion
matrices $C(n),C_Q(n)$. They are exact rational functions in $n$,
not fitted floating-point recurrences. A rational entry is stored as
two lists of coefficients in ascending powers of $n$, for its numerator
and denominator. The replay performs the following checks.

First, with $e_0=(1,0,0,0)^T$ and
$E_n=(8n+1)(9n+1)(10n+1)(11n+a)$, it verifies the rational identities
\begin{equation}\label{eq:paired-transport}
 M(n)K(n+1)Q(n)^T=c(n)K(n),\qquad K(n)u_n^T=-E_ne_0.
\end{equation}
The rank-three reduction is explicit. Form the Krylov matrix
\[
 U(n)=\bigl[e_0,\ M(n)e_0,\ M(n)M(n+1)e_0\bigr].
\]
Then $M(n)U(n+1)=U(n)C(n)$; the analogous $3\times3$ identity
holds for $Q,C_Q$. All denominators, all required determinants, and
the leading recurrence coefficients are nonzero at every integer
$n\geq0$. This is certified by a polynomial sign test: after a
fixed shift the coefficients have one sign and the constant is nonzero,
with the finitely preceding integers checked exactly.

For context, $K$ comes from an Euler Lagrange identity. With
$\partial=z\,d/dz$, use
\[
 \mathcal L=z\prod_{j=1}^5(\partial-a_j+1)
                  -\prod_{j=1}^5(\partial-b_j),
\]
where
\[
 (a_j)=(-16n,-17n,1+11n+\alpha,1+12n+\alpha,2+10n+\alpha),\qquad
 (b_j)=(-2n,-n,\alpha,n+\alpha,2n+\alpha).
\]
If $\mathcal L=\sum_j l_j(z)\partial^j$, then
\[
 \partial\sum_{j\geq1}\sum_{h=0}^{j-1}
       (-1)^h\partial^h(l_jf)\partial^{j-1-h}g
       =f\mathcal Lg-g\mathcal L^*f.
\]
Apply the contiguous change
$f=z^{2n}(\partial+12n+a)F/(12n+a)$ and evaluate at $z=1$.
The fourth $F$ derivative reduces by its hypergeometric equation;
the extra $g$-jet column vanishes. The remaining bilinear matrix is
$K$. The certificate proof only needs~\eqref{eq:paired-transport},
so the supplied rational data can be checked without importing a
computer algebra package for hypergeometric transformations.

Second, the scalar recurrences must describe the actual sums.
For each quarter, \path{problems.py} constructs five explicit polynomial
problems over $\Q[n][k]$ of the form
\begin{equation}\label{eq:paired-certificate}
 \mathcal A(n,k)P(n,k+1)-\mathcal B(n,k)P(n,k)
                                             =\mathcal F(n,k).
\end{equation}
They are the forward scalar recurrence, the dual scalar recurrence,
and the three coordinate transports of the dual row. The right side
is obtained by putting the relevant shifted hypergeometric summands
over their common denominator. Every linear factor, shift and factorial
normalization is retained explicitly in that file. The solver uses
\emph{exact descending coefficients}, then independently substitutes
the full resulting polynomial in~\eqref{eq:paired-certificate}.
For these problems $\deg_k\mathcal A=\deg_k\mathcal B=5$, their
leading coefficients are one and their next coefficients differ by
$-d$, where $d$ is the prescribed degree of $P$. The coefficient
of $k^{i+4}$ has pivot $i-d$. All lower coefficients are consequently
determined by the leading one; the four remaining equations fix that
one by polynomial division. A nonzero residual rejects the certificate.

The degrees $(\deg_nP,\deg_kP)$ are
\[
\begin{array}{c|ccccc}
\alpha&\text{forward}&\text{dual scalar}&\text{row }0&\text{row }1&\text{row }2\\\hline
-3/4&(553,114)&(448,107)&(125,35)&(126,35)&(127,35)\\
-1/4&(553,114)&(446,107)&(124,35)&(125,35)&(126,35).
\end{array}
\]
These large degrees explain the use of machine-readable coefficients.
The complete objects are generated and checked by
\path{checks/paired/verify.py}; no recurrence is accepted because it
matches a finite run of examples.

Third, transfer from a polynomial identity to a sum requires its boundary
terms. For the forward sum the factorial continuation terminates at
$k=0$ and $k=14n+43$ for the four shifted summands. For the derivative
tails use the \emph{uncentered} kernel $\mathcal H_n$ in
\eqref{eq:paired-T}, sampled at $s=m+\alpha$, $m>4n$.
For $n\geq2$, the four shifted tails can all start at $m=12n$:
omitted first and second derivatives lie in their common triple-zero
interval. The lower certificate has a zero of order three there;
its denominator is nonzero by the exact shifted sign test.
At infinity it is constant plus $O(s^{-1})$, so its first and second
derivatives vanish. Absolute locally uniform tail convergence permits
differentiation. The two exceptional recurrence indices $n=0,1$ are
checked exactly for both numerators.

For the dual rows the three certificates become rational identities
\[
 (-s)^j\mathcal R_{n+1}(s)
 -\sum_{i=0}^2Q(n)_{ij}(-s)^i\mathcal R_n(s)
                           =V_j(n,s+1)-V_j(n,s).
\]
At $n\geq1$ all finite poles of $V_j$ lie in $-\alpha+\Z$;
the apparent positive integer poles cancel against numerator zeros.
Translation preserves the total Laurent coefficient of each order,
so $A_n$ and $B_n$ separately transport by $Q(n)$.
For the differentiated sums in~\eqref{eq:paired-state}, both tails
can start at $m=16n$. Omitted terms lie in the common double-zero
interval, the lower certificate derivative vanishes, and its derivative
is $O(s^{-2})$ at infinity. This proves transport of $G_n$;
subtracting the separate $A,B$ transports proves transport of $D$.
The exceptional $n=0$ step checks all nine rational row coordinates.
No assertion of numerical period independence enters this step.

Finally let $\Phi_M(n)=M(0)\cdots M(n-1)$, and similarly for $Q$.
The forward readout rows before multiplication by
$C_n^\alpha=B_n^\alpha(0)$ are
\[
 \rho_q=(1,0,0,0),\quad
 \rho_3=(0,0,-1/2,(2\alpha+1)/4),\quad
 \rho_4=(0,0,0,-1/6).
\]
The preceding recurrences and their first three values identify
$C_n^\alpha\rho_j\Phi_M(n)e_0$ with the forward quantities.
At $n=0$, multiplication by $K(0)$ gives respectively
$-(a/6)(B_0\times A_0)$, $(a/6)(D_0\times A_0)$ and
$-(a/6)(D_0\times B_0)$. Telescope~\eqref{eq:paired-transport}
and use $(vQ)\times(wQ)=\det(Q)(v\times w)Q^{-T}$.
This gives~\eqref{eq:paired-minors} with scalar
\[
 \frac{C_n^\alpha a\prod_{j<n}c(j)}{6E_n\det\Phi_Q(n)}.
\]
Its initial value is $1/6$; the one-step ratio agrees identically
with~\eqref{eq:paired-lambda}, completing the proof. All the stated
matrix, polynomial, boundary-degree and initial-value checks are included
in the standalone replay.

\subsection{Character projection, rates and denominator normalization}
The elementary identity
$\zeta(4,1/4)-\zeta(4,3/4)=256\theta$ gives, with $a=1/4,b=3/4$,
\begin{equation}\label{eq:paired-W}
 W_n=Q_n^aQ_n^b\theta
       -\frac{P_{4,n}^aQ_n^b-P_{4,n}^bQ_n^a}{256}.
\end{equation}
Thus all auxiliary constants disappear. The coefficient of the target
is a product of two different denominators, which matters in estimating
its size.

For reference, we state precisely the analytic input to the arithmetic
comparison below. With
\[
 \sigma=65.45153656\ldots,\qquad g=112.80753924\ldots,
\]
the fixed-quarter construction has
\begin{equation}\label{eq:paired-rates}
 \frac{\log Q_n^a}{n}\longrightarrow\sigma,\qquad
 \left|\zeta(s,a)-\frac{P_{s,n}^a}{Q_n^a}\right|
               \leq e^{-(g-o(1))n},\quad s=3,4.
\end{equation}
Here the analytic recurrence argument is the dual-convergence result
of Leibtag et al.~\cite[Theorem 2.6, Corollary 2.7 and Appendix B.1]{Leibtag},
with the shifted matrices just certified. The required checks can be
made explicitly as follows. Under $\operatorname{diag}(1,n+1,(n+1)^2)$
the two $Q(n)$ limits coincide. Their eigenvalues are a real
$\lambda_0=7.2027\ldots\,10^{-24}$ and a conjugate pair of modulus
$e^{59.5199\ldots}$; $g=\log(|\lambda_1|/\lambda_0)$.
The vector $u_n$ tends, after scaling, to $(1,10,100)$ and has
nonzero projection on the required eigenspace. The exact limiting
matrix and its cyclic-vector determinant are checked in the accompanying
\path{checks/paired/rates.py} replay.

To identify the actual state as the small solution, one does not
assume the conclusion of a numerical recurrence fit. On $2n\leq s\leq17n$,
the products in~\eqref{eq:paired-R}, and their first derivatives,
are bounded by a polynomial in $n$ times $e^{c_0n}$, where
\[
 c_0=14\log15+16\log16-11\log2-11\log3-8\log4
          =51.4744\ldots<\log|\lambda_1|.
\]
For example bound the numerator factors by $15n$ and $16n$, and
the three denominator blocks below by $2n+\alpha$, $3n+\alpha$,
$4n+\alpha$. Fixed shifts change these exponential bounds only by
bounded or polynomial factors. Differentiating products adds only
polynomial factors, including at numerator zeros. For $s\geq17n$,
all numerator factors have magnitude at most $s$, all denominator
factors are at least $s-3/4$, and differentiation gives an integrable
polynomial tail $O(n s^{k-5})$, $k\leq2$. The state therefore has
growth strictly below the conjugate pair. It is nonzero: its initial
coefficient matrix is invertible and its constant coordinate is $1$.
The cited separation argument identifies it with the recessive mode.
The initial readouts and nonzero limiting projection then give the
relative-error estimate in~\eqref{eq:paired-rates}.

The coefficient growth follows independently from the positive sum
\eqref{eq:paired-forward}. Stirling's formula and the Laplace maximum
for $k/n\in[0,14]$ give $\sigma=\max\phi(x)$, where, with $L(t)=t\log t$
and $L(0)=0$,
\[
 \phi(x)=\sum_{j=12}^{14}\bigl(L(x+j)-L(10)-L(x+j-10)\bigr)
       +2L(15)-L(x+1)-L(14-x)-L(x)-L(15-x).
\]
Its second derivative is strictly negative on $(0,14)$, and its first
derivative tends to opposite infinities at the endpoints. Its derivative
therefore vanishes at the unique root in $(0,14)$ of
\[
 \prod_{j\in\{0,1,2,3,4\}}(x+j)
            -\prod_{j=12}^{14}(x+j)(14-x)(15-x)=0.
\]
Fixed quarter shifts do not affect this exponential rate.
Consequently~\eqref{eq:paired-W} has upper error exponent
$2\sigma-g=18.0955\ldots$ before primitive normalization. This is
an upper bound, not an assertion that the projected error actually
grows or has that exact exponent.

A convenient uniform integer denominator is also explicit. Let $O_m$
be the odd part of $\operatorname{lcm}(1,\ldots,m)$ and set, for $n\geq1$,
\begin{equation}\label{eq:paired-clear}
 D_n=6\,2^{60n+3v_2((10n)!)+\lfloor\log_2(16n)\rfloor}O_{76n}^4.
\end{equation}
It clears $Q_n^a,P_{3,n}^a,P_{4,n}^a$ at both quarters, so
$256D_n^2$ clears~\eqref{eq:paired-W}. To verify this, each degree-$m$
quarter-binomial is integral at odd primes by counting multiples in
an arithmetic progression of step four; its exact dyadic cost is
$2m+v_2(m!)$. The derivative identity
\[
 \frac{d}{dx}\binom{x}{m}
       =\sum_{j=1}^m\frac{(-1)^{j-1}}j\binom{x}{m-j}
\]
and the simple-zero tail formulas add one factor
$\operatorname{lcm}(1,\ldots,16n)$. The harmonic sums have index at
most $19n$ and odd denominator dividing $O_{76n}^s$. Combining the
one derivative layer with at most three harmonic layers, and clearing
$1/2,1/3$, proves~\eqref{eq:paired-clear}. This bound is sufficient,
and need not equal the primitive multiplier.

\subsection{What the transferred arithmetic estimates can certify}
It is possible to quantify a limitation without claiming a theorem
about all actual common factors. The following calculation concerns
only the termwise forward valuations and leading-Laurent-coefficient
bounds in the dual rows.
For $p>3$, $p^2>120n$, put $x=n/p$. The two quarter arguments give
phases $t=1/4,3/4$, interchanged when the residue class of $p$ changes.
Define the floor functions
\begin{align*}
 f_t(x,r)={}&\sum_{j=12}^{14}\bigl(
  \lfloor(r+j)x-t\rfloor-\lfloor(r+j-10)x-t\rfloor-\lfloor10x\rfloor\bigr)\\
 &+2\lfloor15x\rfloor-\lfloor(1+r)x\rfloor-\lfloor(14-r)x\rfloor
                         -\lfloor rx\rfloor-\lfloor(15-r)x\rfloor,\\
 \nu_t(x,r)={}&\lfloor(16+r)x-t\rfloor-\lfloor(2+r)x-t\rfloor
              +\lfloor(17+r)x-t\rfloor-\lfloor(1+r)x-t\rfloor\\
 &-\sum_{j\in\{r,11-r,r-1,12-r,r-2,10-r\}}\lfloor jx\rfloor,\\
 \varphi_t(x)={}&\sum_j e_j\lfloor jx\rfloor+
                            \sum_j f_j\lfloor jx-t\rfloor,
\end{align*}
where the two coefficient lists are
\[
 (e_8,e_9,e_{11},e_{12},e_{14},e_{16})=(1,2,2,1,-1,-1),
\]
\[
 (f_{11},f_{12},f_{13},f_{14},f_{16},f_{17},f_{18},f_{19})
                              =(1,2,2,1,-1,-2,-2,-1).
\]
They come respectively from a forward summand at $k=rn$, a triple
pole at $j=rn$, and the scalar~\eqref{eq:paired-lambda}. For the
specified estimates an optimistic ceiling on the quarter's lower
valuation bound for $Q_n^a$ is
\begin{equation}\label{eq:paired-ceiling}
 U_t(x)=\max\left\{\min_{0\leq r\leq14}f_t(x,r),\quad
          \varphi_t(x)+2\min_{3\leq r\leq9}\nu_t(x,r)-1\right\}.
\end{equation}
Indeed, restricting a minimum to selected terms makes it larger.
The two dual rows contributing to $Q$ have lower bounds at most the
selected leading valuation and that valuation minus one, respectively.
Taking the better of the two presentations gives the maximum. Costs
from $P_4$ have been discarded altogether. Every relaxation makes
this a more generous ceiling for the specified estimates.

The functions are periodic with period one and vanish near zero.
All breakpoints on $[0,1]$ belong to
\[
 \left\{\frac{4m+l}{4d}:1\leq d\leq28,\ l\in\{0,1,3\},\
                         0\leq4m+l\leq4d\right\}.
\]
The union has 854 open cells. One rational midpoint gives the exact
constant value on each cell. Integrating through $40$ uses only the
rational weights $1/a-1/b$; the remaining tails are at most $9/40$
and $11/40$, the respective essential maxima divided by $40$.
The standalone calculation \path{checks/paired/arithmetic.py} gives
\begin{equation}\label{eq:paired-density}
 \int_0^\infty\frac{U_{1/4}(x)+U_{3/4}(x)}{x^2}\dd x
                        <88.208.
\end{equation}
To pass from limiting floors to finite prime sums, first restrict $x$
to a bounded set away from the rational walls. The $O(1/p)$ offsets
then preserve all floors; the prime number theorem gives the integral.
Small neighborhoods of the finitely many walls contribute arbitrarily
little because the one-level exponents are uniformly bounded. A bound
of $32$ per quarter suffices: forward selected valuations are at most
$5$, triple-pole leading valuations at most $8$, and the balanced scalar
valuation at most $12$. The tail $p<n/X$ is thus at most $64n/X+o(n)$;
primes below $\sqrt{120n}$ contribute $o(n)$, since each has $O(\log n)$
bounded layers. Let $n$, then $X$, tend to infinity. The sum of the two
phases is independent of $p\bmod4$, so no prime number theorem in
progressions is required for this passage.

For these same estimates the dyadic cost of each quarter's $Q$ is
$90n+O(\log n)$: the three quarter-binomials cost
$60n+3v_2((10n)!)$, and ordinary binomial carries contribute only
$O(\log n)$. The selected dual bound has the same leading cost.
Even after forgiving every $P_4$ cost, success with
\eqref{eq:paired-rates} would require odd-prime savings exceeding
\[
 2\sigma+180\log2-g=142.8620\ldots,
\]
well above~\eqref{eq:paired-density}. This establishes the claimed
limitation of the transferred valuation estimates. It does not bound
the actual gcd of the sums or minors from above, and it does not
exclude faster decay after the character projection.

\paragraph{What remains.}
The paired identities are available, the auxiliary constants cancel
exactly, and both coefficient presentations can be reconstructed.
What is missing is eventual or suitably frequent nonvanishing of
$W_n$, together with shrinking values after its \emph{full} primitive
normalization. The five recorded indices $n=1,2,3,4,6$ give nonzero
but large primitive errors; those computations are finite evidence.
Cancellation inside the sums, a cheaper character projection, or a
stronger projected-error estimate would go beyond the specific
arithmetic comparison just proved.
